\section*{Capítulo \ref{ch:g6:measure} --- Perímetro, área e volume}

\begin{solution}{exo:g6:measure:1}
$3.5$ m $= 350$ cm; \quad $420$ mm $= 42$ cm; \quad
$0.8$ km $= 800$ m; \quad $25\,000$ m $= 25$ km.
\end{solution}

\begin{solution}{exo:g6:measure:2}
Retângulo: $2 \times (8 + 3.5) = 2 \times 11.5 = 23$ cm.
Quadrado: $4 \times 6.2 = 24.8$ cm.
Triângulo: $5 + 7 + 9 = 21$ cm.
\end{solution}

\begin{solution}{exo:g6:measure:3}
Uma volta: $2\pi \times 50 = 100\pi \approx 314$ m. Para $2$ km
$= 2000$ m: $2000 \div 314 \approx 6.4$, então são necessárias $7$
voltas completas ($6$ voltas dão só uns $1\,885$ m).
\end{solution}

\begin{solution}{exo:g6:measure:4}
Retângulo: $9 \times 4 = 36$ cm$^2$. Quadrado: $7^2 = 49$ m$^2$.
\omterm{def:g6:shapes:triangles}{Triângulo retângulo}: $\frac{6 \times 10}{2} = 30$ cm$^2$.
Disco: $\pi \times 3^2 = 9\pi \approx 28$ cm$^2$.
\end{solution}

\begin{solution}{exo:g6:measure:5}
$3$ m$^2 = 30\,000$ cm$^2$; \quad
$45\,000$ cm$^2 = 4.5$ m$^2$; \quad
$2.5$ L $= 2\,500$ cm$^3$; \quad
$4\,500$ L $= 4.5$ m$^3$.
\end{solution}

\begin{solution}{exo:g6:measure:6}
Cerca (\omterm{def:g6:measure:perimeter}{perímetro}): $2 \times (120 + 85) = 2 \times 205 = 410$ m.
\omterm{def:g6:measure:area}{Área}: $120 \times 85 = 10\,200$ m$^2$.
\end{solution}

\begin{solution}{exo:g6:measure:7}
Bloco: $5 \times 4 \times 10 = 200$ cm$^3$. \omterm{def:g5:solids:def}{Cubo}:
$3 \times 3 \times 3 = 27$ cm$^3$.
\end{solution}

\begin{solution}{exo:g6:measure:8}
\omterm{def:g6:measure:perimeter}{Perímetro} $16$ cm significa comprimento $+$ largura $= 8$ cm. Por
exemplo $6 \times 2$ (\omterm{def:g6:measure:area}{área} $12$ cm$^2$) e $5 \times 3$ (\omterm{def:g6:measure:area}{área}
$15$ cm$^2$) --- ou o quadrado $4 \times 4$ (\omterm{def:g6:measure:area}{área} $16$ cm$^2$). Quanto
mais perto de um quadrado, maior a \omterm{def:g6:measure:area}{área}.
\end{solution}

\begin{solution}{exo:g6:measure:9}
\omterm{def:g6:measure:area}{Área}: $10 \times 2 + 2 \times 6 = 20 + 12 = 32$ cm$^2$.

\omterm{def:g6:measure:perimeter}{Perímetro} (haste centrada sob a barra): percorrendo o contorno,
\[
10 + 2 + 4 + 6 + 2 + 6 + 4 + 2 = 36 \text{ cm}
\]
(topo; ponta direita da barra; \omterm{def:g5:solids:def}{face} de baixo à direita; lado direito da
haste; base da haste; lado esquerdo da haste; \omterm{def:g5:solids:def}{face} de baixo à esquerda;
ponta esquerda da barra).
\end{solution}

\begin{solution}{exo:g6:measure:10}
\emph{1.} $V = 25 \times 10 \times 2 = 500$ m$^3$.

\emph{2.} $500$ m$^3 = 500 \times 1000 = 500\,000$ L.

\emph{3.} $500\,000 \div 50\,000 = 10$ horas.
\end{solution}

\begin{solution}{exo:g6:measure:11}
Jardim: $20^2 = 400$ m$^2$. Lago: $\pi \times 3^2 = 9\pi$ m$^2$.
Grama: $400 - 9\pi \approx 400 - 28.3 \approx 372$ m$^2$.
\end{solution}

\begin{solution}{exo:g6:measure:12}
\emph{1.} Barra pequena: $15 \times 6 \times 1 = 90$ cm$^3$. Barra
grande: $30 \times 12 \times 2 = 720$ cm$^3$. O \omterm{def:g6:measure:volume}{volume} fica
multiplicado por $720 \div 90 = 8$ --- dobrar cada uma das três
dimensões multiplica o \omterm{def:g6:measure:volume}{volume} por $2 \times 2 \times 2 = 8$.

\emph{2.} Oito vezes mais chocolate por quatro vezes o preço: sim, a
barra grande vale duas vezes mais a pena por grama.
\end{solution}

\begin{solution}{pb:g6:measure:1}
\textbf{1.} $1 \times 9$: \omterm{def:g6:measure:perimeter}{perímetro} $2 \times (1 + 9) = 20$ m, \omterm{def:g6:measure:area}{área}
$9$ m$^2$. \ $2 \times 8$: \omterm{def:g6:measure:perimeter}{perímetro} $2 \times 10 = 20$ m, \omterm{def:g6:measure:area}{área}
$16$ m$^2$.

\textbf{2.} O \omterm{def:g6:measure:perimeter}{perímetro} é o dobro de (comprimento $+$ largura)
(\cref{ex:g6:measure:perimeter}), então comprimento $+$ largura
$= 20 \div 2 = 10$ m. A tabela:
\begin{center}
\renewcommand{\arraystretch}{1.2}
\begin{tabular}{c|ccccc}
curral & $1 \times 9$ & $2 \times 8$ & $3 \times 7$ & $4 \times 6$
& $5 \times 5$ \\
\hline
\omterm{def:g6:measure:area}{área} (m$^2$) & $9$ & $16$ & $21$ & $24$ & $25$ \\
\end{tabular}
\end{center}
O quadrado $5 \times 5$ é o melhor.

\textbf{3.} $4.5 \times 5.5 = 24.75$ e $4.9 \times 5.1 = 24.99$: cada
vez mais perto de $25$, mas ainda abaixo. Conjectura: o quadrado vence
\emph{todo} retângulo de \omterm{def:g6:measure:perimeter}{perímetro} $20$, lados decimais incluídos.

\textbf{4.} Pela tabela: o \omterm{def:g2:mult:def}{produto} de dois números de \omterm{def:g1:addition:def}{soma} $10$ é maior
quando os dois números são \emph{iguais} ($5$ e $5$). Regra geral: com
\omterm{def:g1:addition:def}{soma} fixa, o \omterm{def:g2:mult:def}{produto} é maior quando as partes são iguais --- quanto
mais desigual o par, menor o \omterm{def:g2:mult:def}{produto}.

\textbf{5.} A partir do quadrado $5 \times 5$, retire a fileira de cima
--- uma faixa de $5$ quadrados unitários ---, deixando um curral
$5 \times 4$; depois cole uma coluna de $4$ quadrados unitários numa
das pontas, transformando o curral em $6 \times 4$. Cinco quadrados
foram tirados e apenas quatro voltaram: a faixa acrescentada fica ao
longo do lado $4$, que é mais curto que o lado $5$ coberto pela faixa
retirada. Perda líquida: um metro quadrado ($25 \to 24$). O passo
seguinte, de $6 \times 4$ para $7 \times 3$, troca uma fileira de $6$
por uma coluna de $3$ e perde mais três ($24 \to 21$): quanto mais
longe do quadrado, mais comprida a faixa que o curral entrega e mais
curta a que ele recebe.

\textbf{6.} Com $w + L + w = 20$, isto é, $L = 20 - 2w$:
\begin{center}
\renewcommand{\arraystretch}{1.2}
\begin{tabular}{c|ccccccc}
$w$ & $1$ & $2$ & $3$ & $4$ & $5$ & $6$ & $7$ \\
\hline
$L$ & $18$ & $16$ & $14$ & $12$ & $10$ & $8$ & $6$ \\
\omterm{def:g6:measure:area}{área} & $18$ & $32$ & $42$ & $48$ & $50$ & $48$ & $42$ \\
\end{tabular}
\end{center}
O vencedor é o $5 \times 10$, de \omterm{def:g6:measure:area}{área} $50$ m$^2$ --- duas vezes mais
comprido (ao longo da parede) do que largo: \emph{não} é um quadrado.

\textbf{7.} Reflita o curral em relação à parede: curral mais imagem
espelhada formam um curral duplicado que usa a cerca duas vezes, $40$ m
de cerca em toda a volta (o lado da parede agora é interno). Pela Parte
I, o melhor retângulo de \omterm{def:g6:measure:perimeter}{perímetro} $40$ é o quadrado $10 \times 10$, de
\omterm{def:g6:measure:area}{área} $100$ m$^2$. O melhor curral de verdade é a \omterm{def:g3:division:half}{metade} desse curral
duplicado: $10$ m ao longo da parede e $5$ m de profundidade, \omterm{def:g6:measure:area}{área}
$50$ m$^2$ --- exatamente o vencedor da tabela, agora explicado.

\textbf{8.} Comprimento $20 = 2 \times \pi \times r$, então
$r = 20 \div 6.28 \approx 3.18$ m. \omterm{def:g6:measure:area}{Área}:
$\pi r^2 \approx 3.14 \times 3.18 \times 3.18 \approx 31.8$ m$^2$ ---
bem mais que os $25$ m$^2$ do quadrado. Dido sabia o que estava
fazendo: para um comprimento dado de contorno, o círculo cerca o
máximo.

\textbf{9.} \omterm{def:g6:measure:perimeter}{Perímetros}: $1 \times 36$: $74$ m; \ $2 \times 18$:
$40$ m; \ $3 \times 12$: $30$ m; \ $4 \times 9$: $26$ m; \
$6 \times 6$: $24$ m. De novo o quadrado --- menos cerca para a \omterm{def:g6:measure:area}{área}
dada. É a Parte I ao contrário: \omterm{def:g6:measure:perimeter}{perímetro} fixo com \omterm{def:g6:measure:area}{área} máxima, ou \omterm{def:g6:measure:area}{área}
fixa com \omterm{def:g6:measure:perimeter}{perímetro} mínimo: as duas coroam o quadrado (e, além de todos
os retângulos, o círculo).

\textbf{10.} Uma parede redonda é o contorno mais curto para o espaço
que ela cerca (questão 8), então uma lata, um cano ou uma caixa d'água
redonda guarda o máximo de conteúdo com o mínimo de material --- e
material é dinheiro.

\textbf{11.} Por exemplo $50$ m $\times$ $0.01$ m: \omterm{def:g6:measure:perimeter}{perímetro}
$2 \times (50 + 0.01) = 100.02$ m, \omterm{def:g6:measure:area}{área} $50 \times 0.01 = 0.5$ m$^2$.
Comprido e fino: quilômetros de contorno e quase nenhuma terra.

\textbf{12.} Falso. O curral da questão 11 tem \omterm{def:g6:measure:perimeter}{perímetro} $100.02$ m e
\omterm{def:g6:measure:area}{área} $0.5$ m$^2$, enquanto o curral $5 \times 5$ tem \omterm{def:g6:measure:perimeter}{perímetro} $20$ m e
\omterm{def:g6:measure:area}{área} $25$ m$^2$: \omterm{def:g6:measure:perimeter}{perímetro} maior, \omterm{def:g6:measure:area}{área} bem menor.

\textbf{13.} \omterm{def:g6:measure:area}{Área}: falta um quadrado unitário, $24 - 1 = 23$ cm$^2$.
\omterm{def:g6:measure:perimeter}{Perímetro}: percorrendo o contorno, o entalhe substitui $1$ cm de parede
reta por três lados do quadradinho, $1 + 1 + 1 = 3$ cm: o \omterm{def:g6:measure:perimeter}{perímetro}
cresce de $20$ cm para $20 - 1 + 3 = 22$ cm. Retirar material
\emph{alongou} o contorno.

\textbf{14.} Cada ampliação sobre um litoral real revela novas baías,
rochas e enseadas --- entalhes sobre entalhes, e a questão 13 mostra
que cada entalhe acrescenta comprimento de contorno sem quase mudar a
\omterm{def:g6:measure:area}{área}. Então o comprimento medido continua crescendo com o nível de
detalhe (o ``paradoxo do litoral''), enquanto a \omterm{def:g6:measure:area}{área} medida se
estabiliza: \omterm{def:g6:measure:perimeter}{perímetro} e \omterm{def:g6:measure:area}{área} são grandezas realmente independentes.

\textbf{15.} Argumento do espelho: o curral duplicado usaria
$2 \times 36 = 72$ m de cerca, e o melhor retângulo de \omterm{def:g6:measure:perimeter}{perímetro} $72$ é
o quadrado $18 \times 18$. Então o melhor curral tem $18$ m ao longo da
parede e $9$ m de profundidade: conferência da cerca,
$9 + 18 + 9 = 36$ m; \omterm{def:g6:measure:area}{área} $18 \times 9 = 162$ m$^2$. Em campo aberto, o
melhor é o quadrado $9 \times 9$, de \omterm{def:g6:measure:area}{área} $81$ m$^2$. A parede do
celeiro exatamente \emph{dobra} a terra do agricultor.
\end{solution}
